\chapter[The CNR model]{The classical normal regression (CNR) model}\label{chap_ClassNormReg}

\section{The CNR model}
\subsection{Basic set-up}
\subsection{The model in matrix notation}

\section{Different approaches to OLS}
\subsection{Least squares estimator}

\subsection{Geometric approach to OLS}
\subsubsection{Projection matrices and their relationships}
Let the matrix $\boldsymbol{P_{X}}$ be defined as

\begin{eqnarray}
\boldsymbol{P_{X}} \equiv \boldsymbol{X}(\boldsymbol{X}'\boldsymbol{X})^{-1}\boldsymbol{X}'.\label{eqn_CNR_Px}
\end{eqnarray}

$\boldsymbol{P_{X}}$ belongs to the general class of so-called projection matrices which have a number of interesting properties. Most importantly, the projection matrix $\boldsymbol{P_{X}}$ is a square matrix of dimension $n \times n$ that gives a vector space projection from $\mathbb{R}^{n}$ to a given subspace. Projection matrices are said to be \emph{idempotent} which implies that $\boldsymbol{P_{X}}\boldsymbol{P_{X}}=\boldsymbol{P_{X}}$.

This important property is straightforward to demonstrate in the case of equation \eqref{eqn_CNR_Px}, where we can use the fact that $$(\boldsymbol{X}'\boldsymbol{X})^{-1}\boldsymbol{X}'\boldsymbol{X}=\boldsymbol{I}$$ in order to show that $$\boldsymbol{P_{X}}\boldsymbol{P_{X}}=\boldsymbol{X}(\boldsymbol{X}'\boldsymbol{X})^{-1}\boldsymbol{X}'\boldsymbol{X}(\boldsymbol{X}'
\boldsymbol{X})^{-1}\boldsymbol{X}'=\boldsymbol{X}(\boldsymbol{X}'\boldsymbol{X})^{-1}\boldsymbol{X}'=\boldsymbol{P_{X}}$$. Furthermore, projection matrices can generally be classified either as \emph{orthogonal} projection matrices or \emph{oblique} (non-orthogonal) projection matrices.

In addition to being square, a projection matrix is also a symmetric matrix if and only if the vector space projection is orthogonal. $\boldsymbol{P_{X}}$ is an orthogonal projection matrix that spans the space $\mathcal{S}(\boldsymbol{X})$; consequently, any point in $\boldsymbol{X}$ that is projected by $\boldsymbol{P_{X}}$ will be projected onto itself. This implies that $\boldsymbol{P_{X}}\boldsymbol{X} = \boldsymbol{X}$. Again, this is easily verified by using equation \eqref{eqn_CNR_Px} which yields $\boldsymbol{P_{X}}\boldsymbol{X}=\boldsymbol{X}(\boldsymbol{X}'\boldsymbol{X})^{-1}\boldsymbol{X}'\boldsymbol{X} = \boldsymbol{X}$.

$\boldsymbol{M_{X}}$ is the orthogonal complement to $\boldsymbol{P_{X}}$, such that $\boldsymbol{M_{X}} = (\boldsymbol{I} - \boldsymbol{P_{X}})$. Consequently,

\begin{eqnarray}
\boldsymbol{P_{X}}\boldsymbol{M_{X}} &=& \boldsymbol{P_{X}}(\boldsymbol{I} - \boldsymbol{P_{X}}) \nonumber \\
&=& \boldsymbol{P_{X}}-\boldsymbol{P_{X}}\boldsymbol{P_{X}} = \boldsymbol{P_{X}}-\boldsymbol{P_{X}}=\boldsymbol{0}.
\end{eqnarray}

Since $\boldsymbol{P_{X}}$ and $\boldsymbol{M_{X}}$ are orthogonal to each other, any vector in $\mathcal{S}(\boldsymbol{X})$ will either be projected onto itself (i.e. $\boldsymbol{P_{X}}\boldsymbol{X} = \boldsymbol{X}$) or it will be annihilated ($\boldsymbol{M_{X}}\boldsymbol{X} = \boldsymbol{0}$).

Now, suppose that the matrix $\boldsymbol{X}$ can be decomposed into the submatrices $\boldsymbol{X}_{1}$ and $\boldsymbol{X}_{2}$, such that $\boldsymbol{X} = [\boldsymbol{X}_{1}\hspace{1ex}\boldsymbol{X}_{2}]$. Analogous to $\boldsymbol{P_{X}}$, we can now define the projection matrix $\boldsymbol{P}_{1} = \boldsymbol{X}_{1}(\boldsymbol{X}_{1}'\boldsymbol{X}_{1})^{-1}\boldsymbol{X}_{1}'$. This yields the following important relationships:

\begin{eqnarray}
\boldsymbol{P_{X}}\boldsymbol{P}_{1}=\boldsymbol{P}_{1} \qquad &\text{since}& \qquad \boldsymbol{P}_{1}<\boldsymbol{P_{X}}\\
\boldsymbol{M_{X}}\boldsymbol{M}_{1}=\boldsymbol{M_{X}} \qquad &\text{since}& \qquad \boldsymbol{M}_{1}>\boldsymbol{M_{X}}
\end{eqnarray}


\subsubsection{Geometry of the Frisch-Waugh-Lovell theorem}
See \citet{Lovell2008} for a simple proof that does not rely on matrix algebra. More comprehensive treatment of the Frisch-Waugh-Lovell theorem (FWL) can be found in \citet{Davidson2004}.

\subsection{General Method of Moments (GMM)}

\subsection{Maximum Likelihood Estimation and OLS}
The maximum likelihood estimate (MLE) is that assumed value of an unknown population parameter which maximises the probability of obtaining the sample data that is actually observed.

The normal regression model is the linear regression model under the restriction that the error term $e_{i}$ is independent of $\boldsymbol{x}_{i}$ and has the distribution $(0, \sigma^{2})$ Thus the standard set-up of OLS that we have just seen can easily be re-written as

\begin{eqnarray}
e_{i} | \boldsymbol{x}_{i} \sim (0, \sigma^{2})
\end{eqnarray}

Recall from section~\ref{chap_Review:Stats} that the mean and variance of a variable $V = \mu_{0} + e$ where $e \sim (0,\sigma^{2})$, is $\mu_{V} = \mu$ and $\sigma^{2}_{V}=\sigma^{2}_{e}$. This implies that

\begin{eqnarray}
y_{i} | \boldsymbol{x}_{i} \sim (\boldsymbol{x}_{i}'\boldsymbol{\beta}, \sigma^{2})
\end{eqnarray}

Equivalently, we can rewrite the linear model as $\boldsymbol{y} \sim N{\boldsymbol{\Omega}, \boldsymbol{\Omega}}$ where $\boldsymbol{\Omega}$ is assumed to meet the Gauss-Markov assumptions.

Given the parametric nature of the CNR, we can use likelihood methods\index{likelihood methods} for estimating the parameters. The likelihood function for each observation is

\begin{eqnarray}
\mathcal{L}_{i}  = f(y_{i}) = \frac{1}{\sqrt{2 \pi \sigma^{2}}}e^{-\frac{(y_{i} - \boldsymbol{x}'_{i}\boldsymbol{\beta})^{2}}{2\sigma^{2}}}
\end{eqnarray}

The maximum likelihood estimator\index{maximum likelihood estimator (MLE)|see{likelihood methods}} (MLE; $\hat{\boldsymbol{\beta}}_{mle}, \sigma^{2}_{mle}$) maximises $\log \mathcal{L}(\boldsymbol{\beta}, \sigma^{2})$. Since this is a function of $\boldsymbol{\beta}$ only through the sum of squared errors $SSE_{n}(\boldsymbol{\beta})$, maximising the likelihood is equivalent to minimising the sum of squared errors. Therefore

\begin{eqnarray}
\hat{\boldsymbol{\beta}}_{mle} = \boldsymbol{\hat{\beta}}_{ols}.
\end{eqnarray}


\begin{subappendices}
\section{Code for Spatial Stats in R}
\lstinputlisting{"Programs/UP504_OLS.R"}

\renewcommand\bibname{Further reading}

\bibliographystyle{econometrica}
\bibliography{Literature/OverLord}
\IfFileExists{Literature/ClassNormReg.bbl}{\chapter[The CNR model]{The classical normal regression (CNR) model}\label{chap_ClassNormReg}

\section{The CNR model}
\subsection{Basic set-up}
\subsection{The model in matrix notation}

\section{Different approaches to OLS}
\subsection{Least squares estimator}

\subsection{Geometric approach to OLS}
\subsubsection{Projection matrices and their relationships}
Let the matrix $\boldsymbol{P_{X}}$ be defined as

\begin{eqnarray}
\boldsymbol{P_{X}} \equiv \boldsymbol{X}(\boldsymbol{X}'\boldsymbol{X})^{-1}\boldsymbol{X}'.\label{eqn_CNR_Px}
\end{eqnarray}

$\boldsymbol{P_{X}}$ belongs to the general class of so-called projection matrices which have a number of interesting properties. Most importantly, the projection matrix $\boldsymbol{P_{X}}$ is a square matrix of dimension $n \times n$ that gives a vector space projection from $\mathbb{R}^{n}$ to a given subspace. Projection matrices are said to be \emph{idempotent} which implies that $\boldsymbol{P_{X}}\boldsymbol{P_{X}}=\boldsymbol{P_{X}}$.

This important property is straightforward to demonstrate in the case of equation \eqref{eqn_CNR_Px}, where we can use the fact that $$(\boldsymbol{X}'\boldsymbol{X})^{-1}\boldsymbol{X}'\boldsymbol{X}=\boldsymbol{I}$$ in order to show that $$\boldsymbol{P_{X}}\boldsymbol{P_{X}}=\boldsymbol{X}(\boldsymbol{X}'\boldsymbol{X})^{-1}\boldsymbol{X}'\boldsymbol{X}(\boldsymbol{X}'
\boldsymbol{X})^{-1}\boldsymbol{X}'=\boldsymbol{X}(\boldsymbol{X}'\boldsymbol{X})^{-1}\boldsymbol{X}'=\boldsymbol{P_{X}}$$. Furthermore, projection matrices can generally be classified either as \emph{orthogonal} projection matrices or \emph{oblique} (non-orthogonal) projection matrices.

In addition to being square, a projection matrix is also a symmetric matrix if and only if the vector space projection is orthogonal. $\boldsymbol{P_{X}}$ is an orthogonal projection matrix that spans the space $\mathcal{S}(\boldsymbol{X})$; consequently, any point in $\boldsymbol{X}$ that is projected by $\boldsymbol{P_{X}}$ will be projected onto itself. This implies that $\boldsymbol{P_{X}}\boldsymbol{X} = \boldsymbol{X}$. Again, this is easily verified by using equation \eqref{eqn_CNR_Px} which yields $\boldsymbol{P_{X}}\boldsymbol{X}=\boldsymbol{X}(\boldsymbol{X}'\boldsymbol{X})^{-1}\boldsymbol{X}'\boldsymbol{X} = \boldsymbol{X}$.

$\boldsymbol{M_{X}}$ is the orthogonal complement to $\boldsymbol{P_{X}}$, such that $\boldsymbol{M_{X}} = (\boldsymbol{I} - \boldsymbol{P_{X}})$. Consequently,

\begin{eqnarray}
\boldsymbol{P_{X}}\boldsymbol{M_{X}} &=& \boldsymbol{P_{X}}(\boldsymbol{I} - \boldsymbol{P_{X}}) \nonumber \\
&=& \boldsymbol{P_{X}}-\boldsymbol{P_{X}}\boldsymbol{P_{X}} = \boldsymbol{P_{X}}-\boldsymbol{P_{X}}=\boldsymbol{0}.
\end{eqnarray}

Since $\boldsymbol{P_{X}}$ and $\boldsymbol{M_{X}}$ are orthogonal to each other, any vector in $\mathcal{S}(\boldsymbol{X})$ will either be projected onto itself (i.e. $\boldsymbol{P_{X}}\boldsymbol{X} = \boldsymbol{X}$) or it will be annihilated ($\boldsymbol{M_{X}}\boldsymbol{X} = \boldsymbol{0}$).

Now, suppose that the matrix $\boldsymbol{X}$ can be decomposed into the submatrices $\boldsymbol{X}_{1}$ and $\boldsymbol{X}_{2}$, such that $\boldsymbol{X} = [\boldsymbol{X}_{1}\hspace{1ex}\boldsymbol{X}_{2}]$. Analogous to $\boldsymbol{P_{X}}$, we can now define the projection matrix $\boldsymbol{P}_{1} = \boldsymbol{X}_{1}(\boldsymbol{X}_{1}'\boldsymbol{X}_{1})^{-1}\boldsymbol{X}_{1}'$. This yields the following important relationships:

\begin{eqnarray}
\boldsymbol{P_{X}}\boldsymbol{P}_{1}=\boldsymbol{P}_{1} \qquad &\text{since}& \qquad \boldsymbol{P}_{1}<\boldsymbol{P_{X}}\\
\boldsymbol{M_{X}}\boldsymbol{M}_{1}=\boldsymbol{M_{X}} \qquad &\text{since}& \qquad \boldsymbol{M}_{1}>\boldsymbol{M_{X}}
\end{eqnarray}


\subsubsection{Geometry of the Frisch-Waugh-Lovell theorem}
See \citet{Lovell2008} for a simple proof that does not rely on matrix algebra. More comprehensive treatment of the Frisch-Waugh-Lovell theorem (FWL) can be found in \citet{Davidson2004}.

\subsection{General Method of Moments (GMM)}

\subsection{Maximum Likelihood Estimation and OLS}
The maximum likelihood estimate (MLE) is that assumed value of an unknown population parameter which maximises the probability of obtaining the sample data that is actually observed.

The normal regression model is the linear regression model under the restriction that the error term $e_{i}$ is independent of $\boldsymbol{x}_{i}$ and has the distribution $(0, \sigma^{2})$ Thus the standard set-up of OLS that we have just seen can easily be re-written as

\begin{eqnarray}
e_{i} | \boldsymbol{x}_{i} \sim (0, \sigma^{2})
\end{eqnarray}

Recall from section~\ref{chap_Review:Stats} that the mean and variance of a variable $V = \mu_{0} + e$ where $e \sim (0,\sigma^{2})$, is $\mu_{V} = \mu$ and $\sigma^{2}_{V}=\sigma^{2}_{e}$. This implies that

\begin{eqnarray}
y_{i} | \boldsymbol{x}_{i} \sim (\boldsymbol{x}_{i}'\boldsymbol{\beta}, \sigma^{2})
\end{eqnarray}

Equivalently, we can rewrite the linear model as $\boldsymbol{y} \sim N{\boldsymbol{\Omega}, \boldsymbol{\Omega}}$ where $\boldsymbol{\Omega}$ is assumed to meet the Gauss-Markov assumptions.

Given the parametric nature of the CNR, we can use likelihood methods\index{likelihood methods} for estimating the parameters. The likelihood function for each observation is

\begin{eqnarray}
\mathcal{L}_{i}  = f(y_{i}) = \frac{1}{\sqrt{2 \pi \sigma^{2}}}e^{-\frac{(y_{i} - \boldsymbol{x}'_{i}\boldsymbol{\beta})^{2}}{2\sigma^{2}}}
\end{eqnarray}

The maximum likelihood estimator\index{maximum likelihood estimator (MLE)|see{likelihood methods}} (MLE; $\hat{\boldsymbol{\beta}}_{mle}, \sigma^{2}_{mle}$) maximises $\log \mathcal{L}(\boldsymbol{\beta}, \sigma^{2})$. Since this is a function of $\boldsymbol{\beta}$ only through the sum of squared errors $SSE_{n}(\boldsymbol{\beta})$, maximising the likelihood is equivalent to minimising the sum of squared errors. Therefore

\begin{eqnarray}
\hat{\boldsymbol{\beta}}_{mle} = \boldsymbol{\hat{\beta}}_{ols}.
\end{eqnarray}


\begin{subappendices}
\section{Code for Spatial Stats in R}
\lstinputlisting{"Programs/UP504_OLS.R"}

\renewcommand\bibname{Further reading}

\bibliographystyle{econometrica}
\bibliography{Literature/OverLord}
\IfFileExists{Literature/ClassNormReg.bbl}{\chapter[The CNR model]{The classical normal regression (CNR) model}\label{chap_ClassNormReg}

\section{The CNR model}
\subsection{Basic set-up}
\subsection{The model in matrix notation}

\section{Different approaches to OLS}
\subsection{Least squares estimator}

\subsection{Geometric approach to OLS}
\subsubsection{Projection matrices and their relationships}
Let the matrix $\boldsymbol{P_{X}}$ be defined as

\begin{eqnarray}
\boldsymbol{P_{X}} \equiv \boldsymbol{X}(\boldsymbol{X}'\boldsymbol{X})^{-1}\boldsymbol{X}'.\label{eqn_CNR_Px}
\end{eqnarray}

$\boldsymbol{P_{X}}$ belongs to the general class of so-called projection matrices which have a number of interesting properties. Most importantly, the projection matrix $\boldsymbol{P_{X}}$ is a square matrix of dimension $n \times n$ that gives a vector space projection from $\mathbb{R}^{n}$ to a given subspace. Projection matrices are said to be \emph{idempotent} which implies that $\boldsymbol{P_{X}}\boldsymbol{P_{X}}=\boldsymbol{P_{X}}$.

This important property is straightforward to demonstrate in the case of equation \eqref{eqn_CNR_Px}, where we can use the fact that $$(\boldsymbol{X}'\boldsymbol{X})^{-1}\boldsymbol{X}'\boldsymbol{X}=\boldsymbol{I}$$ in order to show that $$\boldsymbol{P_{X}}\boldsymbol{P_{X}}=\boldsymbol{X}(\boldsymbol{X}'\boldsymbol{X})^{-1}\boldsymbol{X}'\boldsymbol{X}(\boldsymbol{X}'
\boldsymbol{X})^{-1}\boldsymbol{X}'=\boldsymbol{X}(\boldsymbol{X}'\boldsymbol{X})^{-1}\boldsymbol{X}'=\boldsymbol{P_{X}}$$. Furthermore, projection matrices can generally be classified either as \emph{orthogonal} projection matrices or \emph{oblique} (non-orthogonal) projection matrices.

In addition to being square, a projection matrix is also a symmetric matrix if and only if the vector space projection is orthogonal. $\boldsymbol{P_{X}}$ is an orthogonal projection matrix that spans the space $\mathcal{S}(\boldsymbol{X})$; consequently, any point in $\boldsymbol{X}$ that is projected by $\boldsymbol{P_{X}}$ will be projected onto itself. This implies that $\boldsymbol{P_{X}}\boldsymbol{X} = \boldsymbol{X}$. Again, this is easily verified by using equation \eqref{eqn_CNR_Px} which yields $\boldsymbol{P_{X}}\boldsymbol{X}=\boldsymbol{X}(\boldsymbol{X}'\boldsymbol{X})^{-1}\boldsymbol{X}'\boldsymbol{X} = \boldsymbol{X}$.

$\boldsymbol{M_{X}}$ is the orthogonal complement to $\boldsymbol{P_{X}}$, such that $\boldsymbol{M_{X}} = (\boldsymbol{I} - \boldsymbol{P_{X}})$. Consequently,

\begin{eqnarray}
\boldsymbol{P_{X}}\boldsymbol{M_{X}} &=& \boldsymbol{P_{X}}(\boldsymbol{I} - \boldsymbol{P_{X}}) \nonumber \\
&=& \boldsymbol{P_{X}}-\boldsymbol{P_{X}}\boldsymbol{P_{X}} = \boldsymbol{P_{X}}-\boldsymbol{P_{X}}=\boldsymbol{0}.
\end{eqnarray}

Since $\boldsymbol{P_{X}}$ and $\boldsymbol{M_{X}}$ are orthogonal to each other, any vector in $\mathcal{S}(\boldsymbol{X})$ will either be projected onto itself (i.e. $\boldsymbol{P_{X}}\boldsymbol{X} = \boldsymbol{X}$) or it will be annihilated ($\boldsymbol{M_{X}}\boldsymbol{X} = \boldsymbol{0}$).

Now, suppose that the matrix $\boldsymbol{X}$ can be decomposed into the submatrices $\boldsymbol{X}_{1}$ and $\boldsymbol{X}_{2}$, such that $\boldsymbol{X} = [\boldsymbol{X}_{1}\hspace{1ex}\boldsymbol{X}_{2}]$. Analogous to $\boldsymbol{P_{X}}$, we can now define the projection matrix $\boldsymbol{P}_{1} = \boldsymbol{X}_{1}(\boldsymbol{X}_{1}'\boldsymbol{X}_{1})^{-1}\boldsymbol{X}_{1}'$. This yields the following important relationships:

\begin{eqnarray}
\boldsymbol{P_{X}}\boldsymbol{P}_{1}=\boldsymbol{P}_{1} \qquad &\text{since}& \qquad \boldsymbol{P}_{1}<\boldsymbol{P_{X}}\\
\boldsymbol{M_{X}}\boldsymbol{M}_{1}=\boldsymbol{M_{X}} \qquad &\text{since}& \qquad \boldsymbol{M}_{1}>\boldsymbol{M_{X}}
\end{eqnarray}


\subsubsection{Geometry of the Frisch-Waugh-Lovell theorem}
See \citet{Lovell2008} for a simple proof that does not rely on matrix algebra. More comprehensive treatment of the Frisch-Waugh-Lovell theorem (FWL) can be found in \citet{Davidson2004}.

\subsection{General Method of Moments (GMM)}

\subsection{Maximum Likelihood Estimation and OLS}
The maximum likelihood estimate (MLE) is that assumed value of an unknown population parameter which maximises the probability of obtaining the sample data that is actually observed.

The normal regression model is the linear regression model under the restriction that the error term $e_{i}$ is independent of $\boldsymbol{x}_{i}$ and has the distribution $(0, \sigma^{2})$ Thus the standard set-up of OLS that we have just seen can easily be re-written as

\begin{eqnarray}
e_{i} | \boldsymbol{x}_{i} \sim (0, \sigma^{2})
\end{eqnarray}

Recall from section~\ref{chap_Review:Stats} that the mean and variance of a variable $V = \mu_{0} + e$ where $e \sim (0,\sigma^{2})$, is $\mu_{V} = \mu$ and $\sigma^{2}_{V}=\sigma^{2}_{e}$. This implies that

\begin{eqnarray}
y_{i} | \boldsymbol{x}_{i} \sim (\boldsymbol{x}_{i}'\boldsymbol{\beta}, \sigma^{2})
\end{eqnarray}

Equivalently, we can rewrite the linear model as $\boldsymbol{y} \sim N{\boldsymbol{\Omega}, \boldsymbol{\Omega}}$ where $\boldsymbol{\Omega}$ is assumed to meet the Gauss-Markov assumptions.

Given the parametric nature of the CNR, we can use likelihood methods\index{likelihood methods} for estimating the parameters. The likelihood function for each observation is

\begin{eqnarray}
\mathcal{L}_{i}  = f(y_{i}) = \frac{1}{\sqrt{2 \pi \sigma^{2}}}e^{-\frac{(y_{i} - \boldsymbol{x}'_{i}\boldsymbol{\beta})^{2}}{2\sigma^{2}}}
\end{eqnarray}

The maximum likelihood estimator\index{maximum likelihood estimator (MLE)|see{likelihood methods}} (MLE; $\hat{\boldsymbol{\beta}}_{mle}, \sigma^{2}_{mle}$) maximises $\log \mathcal{L}(\boldsymbol{\beta}, \sigma^{2})$. Since this is a function of $\boldsymbol{\beta}$ only through the sum of squared errors $SSE_{n}(\boldsymbol{\beta})$, maximising the likelihood is equivalent to minimising the sum of squared errors. Therefore

\begin{eqnarray}
\hat{\boldsymbol{\beta}}_{mle} = \boldsymbol{\hat{\beta}}_{ols}.
\end{eqnarray}


\begin{subappendices}
\section{Code for Spatial Stats in R}
\lstinputlisting{"Programs/UP504_OLS.R"}

\renewcommand\bibname{Further reading}

\bibliographystyle{econometrica}
\bibliography{Literature/OverLord}
\IfFileExists{Literature/ClassNormReg.bbl}{\chapter[The CNR model]{The classical normal regression (CNR) model}\label{chap_ClassNormReg}

\section{The CNR model}
\subsection{Basic set-up}
\subsection{The model in matrix notation}

\section{Different approaches to OLS}
\subsection{Least squares estimator}

\subsection{Geometric approach to OLS}
\subsubsection{Projection matrices and their relationships}
Let the matrix $\boldsymbol{P_{X}}$ be defined as

\begin{eqnarray}
\boldsymbol{P_{X}} \equiv \boldsymbol{X}(\boldsymbol{X}'\boldsymbol{X})^{-1}\boldsymbol{X}'.\label{eqn_CNR_Px}
\end{eqnarray}

$\boldsymbol{P_{X}}$ belongs to the general class of so-called projection matrices which have a number of interesting properties. Most importantly, the projection matrix $\boldsymbol{P_{X}}$ is a square matrix of dimension $n \times n$ that gives a vector space projection from $\mathbb{R}^{n}$ to a given subspace. Projection matrices are said to be \emph{idempotent} which implies that $\boldsymbol{P_{X}}\boldsymbol{P_{X}}=\boldsymbol{P_{X}}$.

This important property is straightforward to demonstrate in the case of equation \eqref{eqn_CNR_Px}, where we can use the fact that $$(\boldsymbol{X}'\boldsymbol{X})^{-1}\boldsymbol{X}'\boldsymbol{X}=\boldsymbol{I}$$ in order to show that $$\boldsymbol{P_{X}}\boldsymbol{P_{X}}=\boldsymbol{X}(\boldsymbol{X}'\boldsymbol{X})^{-1}\boldsymbol{X}'\boldsymbol{X}(\boldsymbol{X}'
\boldsymbol{X})^{-1}\boldsymbol{X}'=\boldsymbol{X}(\boldsymbol{X}'\boldsymbol{X})^{-1}\boldsymbol{X}'=\boldsymbol{P_{X}}$$. Furthermore, projection matrices can generally be classified either as \emph{orthogonal} projection matrices or \emph{oblique} (non-orthogonal) projection matrices.

In addition to being square, a projection matrix is also a symmetric matrix if and only if the vector space projection is orthogonal. $\boldsymbol{P_{X}}$ is an orthogonal projection matrix that spans the space $\mathcal{S}(\boldsymbol{X})$; consequently, any point in $\boldsymbol{X}$ that is projected by $\boldsymbol{P_{X}}$ will be projected onto itself. This implies that $\boldsymbol{P_{X}}\boldsymbol{X} = \boldsymbol{X}$. Again, this is easily verified by using equation \eqref{eqn_CNR_Px} which yields $\boldsymbol{P_{X}}\boldsymbol{X}=\boldsymbol{X}(\boldsymbol{X}'\boldsymbol{X})^{-1}\boldsymbol{X}'\boldsymbol{X} = \boldsymbol{X}$.

$\boldsymbol{M_{X}}$ is the orthogonal complement to $\boldsymbol{P_{X}}$, such that $\boldsymbol{M_{X}} = (\boldsymbol{I} - \boldsymbol{P_{X}})$. Consequently,

\begin{eqnarray}
\boldsymbol{P_{X}}\boldsymbol{M_{X}} &=& \boldsymbol{P_{X}}(\boldsymbol{I} - \boldsymbol{P_{X}}) \nonumber \\
&=& \boldsymbol{P_{X}}-\boldsymbol{P_{X}}\boldsymbol{P_{X}} = \boldsymbol{P_{X}}-\boldsymbol{P_{X}}=\boldsymbol{0}.
\end{eqnarray}

Since $\boldsymbol{P_{X}}$ and $\boldsymbol{M_{X}}$ are orthogonal to each other, any vector in $\mathcal{S}(\boldsymbol{X})$ will either be projected onto itself (i.e. $\boldsymbol{P_{X}}\boldsymbol{X} = \boldsymbol{X}$) or it will be annihilated ($\boldsymbol{M_{X}}\boldsymbol{X} = \boldsymbol{0}$).

Now, suppose that the matrix $\boldsymbol{X}$ can be decomposed into the submatrices $\boldsymbol{X}_{1}$ and $\boldsymbol{X}_{2}$, such that $\boldsymbol{X} = [\boldsymbol{X}_{1}\hspace{1ex}\boldsymbol{X}_{2}]$. Analogous to $\boldsymbol{P_{X}}$, we can now define the projection matrix $\boldsymbol{P}_{1} = \boldsymbol{X}_{1}(\boldsymbol{X}_{1}'\boldsymbol{X}_{1})^{-1}\boldsymbol{X}_{1}'$. This yields the following important relationships:

\begin{eqnarray}
\boldsymbol{P_{X}}\boldsymbol{P}_{1}=\boldsymbol{P}_{1} \qquad &\text{since}& \qquad \boldsymbol{P}_{1}<\boldsymbol{P_{X}}\\
\boldsymbol{M_{X}}\boldsymbol{M}_{1}=\boldsymbol{M_{X}} \qquad &\text{since}& \qquad \boldsymbol{M}_{1}>\boldsymbol{M_{X}}
\end{eqnarray}


\subsubsection{Geometry of the Frisch-Waugh-Lovell theorem}
See \citet{Lovell2008} for a simple proof that does not rely on matrix algebra. More comprehensive treatment of the Frisch-Waugh-Lovell theorem (FWL) can be found in \citet{Davidson2004}.

\subsection{General Method of Moments (GMM)}

\subsection{Maximum Likelihood Estimation and OLS}
The maximum likelihood estimate (MLE) is that assumed value of an unknown population parameter which maximises the probability of obtaining the sample data that is actually observed.

The normal regression model is the linear regression model under the restriction that the error term $e_{i}$ is independent of $\boldsymbol{x}_{i}$ and has the distribution $(0, \sigma^{2})$ Thus the standard set-up of OLS that we have just seen can easily be re-written as

\begin{eqnarray}
e_{i} | \boldsymbol{x}_{i} \sim (0, \sigma^{2})
\end{eqnarray}

Recall from section~\ref{chap_Review:Stats} that the mean and variance of a variable $V = \mu_{0} + e$ where $e \sim (0,\sigma^{2})$, is $\mu_{V} = \mu$ and $\sigma^{2}_{V}=\sigma^{2}_{e}$. This implies that

\begin{eqnarray}
y_{i} | \boldsymbol{x}_{i} \sim (\boldsymbol{x}_{i}'\boldsymbol{\beta}, \sigma^{2})
\end{eqnarray}

Equivalently, we can rewrite the linear model as $\boldsymbol{y} \sim N{\boldsymbol{\Omega}, \boldsymbol{\Omega}}$ where $\boldsymbol{\Omega}$ is assumed to meet the Gauss-Markov assumptions.

Given the parametric nature of the CNR, we can use likelihood methods\index{likelihood methods} for estimating the parameters. The likelihood function for each observation is

\begin{eqnarray}
\mathcal{L}_{i}  = f(y_{i}) = \frac{1}{\sqrt{2 \pi \sigma^{2}}}e^{-\frac{(y_{i} - \boldsymbol{x}'_{i}\boldsymbol{\beta})^{2}}{2\sigma^{2}}}
\end{eqnarray}

The maximum likelihood estimator\index{maximum likelihood estimator (MLE)|see{likelihood methods}} (MLE; $\hat{\boldsymbol{\beta}}_{mle}, \sigma^{2}_{mle}$) maximises $\log \mathcal{L}(\boldsymbol{\beta}, \sigma^{2})$. Since this is a function of $\boldsymbol{\beta}$ only through the sum of squared errors $SSE_{n}(\boldsymbol{\beta})$, maximising the likelihood is equivalent to minimising the sum of squared errors. Therefore

\begin{eqnarray}
\hat{\boldsymbol{\beta}}_{mle} = \boldsymbol{\hat{\beta}}_{ols}.
\end{eqnarray}


\begin{subappendices}
\section{Code for Spatial Stats in R}
\lstinputlisting{"Programs/UP504_OLS.R"}

\renewcommand\bibname{Further reading}

\bibliographystyle{econometrica}
\bibliography{Literature/OverLord}
\IfFileExists{Literature/ClassNormReg.bbl}{\input{Literature/ClassNormReg.bbl}}{\typeout{Need to run bibtex}}
\end{subappendices}
}{\typeout{Need to run bibtex}}
\end{subappendices}
}{\typeout{Need to run bibtex}}
\end{subappendices}
}{\typeout{Need to run bibtex}}
\end{subappendices}
